\subsection{Spectrum relative to a subalgebra}

In this section, we assume that K = C.

Lemma 3. --- Let \(X_{1}\) and \(X_{2}\) be compact subsets of C. If \(X_{2}\) is contained in \(X_{1}\) and if the boundary of \(X_{1}\) in C is contained in \(X_{2}\) , then \(X_{1}\) is the union of \(X_{2}\) and of certain bounded connected components of the complement of \(X_{2}\) in C.

Let U be a connected component of \(\mathbf{C}-X_{2}\) . Every boundary point of \(X_{1} \cap U\) in the open set U is also a boundary point of \(X_{1}\) in C, hence belongs to \(X_{2}\) by hypothesis; since \(U \cap X_{2} = \varnothing\) , we see that \(X_{1} \cap U\) has no boundary point in the space U. Since U is connected, the intersection \(X_{1} \cap U\) is either empty or equal to U (TG, I, p.~82, cor.), and the lemma follows.

PROPOSITION 6. — Let A be a unital Banach algebra and B a closed unital subalgebra of A. For every \(x \in B\) , one has \(\mathrm{Sp}_{\mathrm{B}}(x) \supset \mathrm{Sp}_{\mathrm{A}}(x)\) , and the boundary of \(\mathrm{Sp}_{\mathrm{A}}(x)\) in \(\mathbf{C}\) contains the boundary of \(\mathrm{Sp}_{\mathrm{B}}(x)\) in \(\mathbf{C}\) . In particular, if \(\mathrm{Sp}_{\mathrm{B}}(x) \subset \mathbf{R}\) , then we have \(\mathrm{Sp}_{\mathrm{B}}(x) = \mathrm{Sp}_{\mathrm{A}}(x)\) .

We have \(\mathrm{Sp}_{\mathrm{B}}(x) \supset \mathrm{Sp}_{\mathrm{A}}(x)\) (remark 6 of I, p.~3). If \(\lambda\) is a point on the boundary of \(\mathrm{Sp}_{\mathrm{B}}(x)\) in C, there exists a sequence \((\lambda_{n})\) of points exterior to \(\mathrm{Sp}_{\mathrm{B}}(x)\) tending to \(\lambda\) . Then \(x - \lambda_{n}\) is invertible in B and tends to \(x - \lambda\) , which is not invertible in B; hence \(x - \lambda\) is a left or right topological zero divisor in B (prop. 4 of I, p.~24), hence in A. Thus, \(\lambda \in \mathrm{Sp}_{\mathrm{A}}(x)\) . But since \(\mathrm{Sp}_{\mathrm{A}}(x) \subset \mathrm{Sp}_{\mathrm{B}}(x)\) , the complex number \(\lambda \in \mathrm{Fr}_{\mathbf{C}}(\mathrm{Sp}_{\mathrm{B}}(x))\) cannot be interior to \(\mathrm{Sp}_{\mathrm{A}}(x)\) , hence belongs to its boundary.

COROLLARY. --- The set \(\mathrm{Sp}_{\mathrm{B}}(x)\) is the union of \(\mathrm{Sp}_{\mathrm{A}}(x)\) and of certain bounded connected components of \(\mathbf{C}-\mathrm{Sp}_{\mathrm{A}}(x)\) .

This follows from prop. 6 and lemma 3.

This corollary will be completed by propositions 13 of I, p.~46 and 14 of I, p.~46.

